% \iffalse meta-comment
%
% Copyright (C) 2017- by Yanshuo Chu <yanshuoc@gmail.com>
%
% This file may be distributed and/or modified under the
% conditions of the LaTeX Project Public License, either version 1.3a
% of this license or (at your option) any later version.
% The latest version of this license is in:
%
% http://www.latex-project.org/lppl.txt
%
% and version 1.3a or later is part of all distributions of LaTeX
% version 2004/10/01 or later.
%
% \fi
%
% \section{学位论文类的术语与符号表}
%
% ^^A ======================================================================
% ^^A Module thesis-glossary: 术语与符号表（glossaries 配置 + 示例条目）（hit-thesis）
% ^^A ======================================================================
%    \begin{macrocode}
%<*thesis-glossary>
%    \end{macrocode}
%
% \changes{v3.1d}{2025/03/03}{将 \cs{glossarystyle} 替换为 \cs{setglossarystyle}。\issue[200]}
% \changes{v3.2a}{2026/08/12}{book-glossary 的两处 english 判断改用 \cs{legacy\_if:nTF}。
%   后面的 \cs{newacronym} 条目含英文字面文本，留在 expl3 之外}
% 在 \pkg{glossaries} 中表示 \cs{glossarystle} 这个命令 Deprecated in v3.08a. Removed in v4.50.
%    \begin{macrocode}
\bool_if:NTF \g__hit_en_bool
  {
    \RequirePackage [ xindy , order=letter ] { glossaries-extra }
%    \end{macrocode}
%
% 样式代码存进 \pkg{glossaries} 时就已经记号化完了，写在里面的
% \cs{ExplSyntaxOn} 不起作用。环境正文本来就是排版内容，直接用 2e 写法。
%    \begin{macrocode}
\newglossarystyle{engstyle}{%
  \setglossarystyle{long}%
  \RenewDocumentEnvironment{theglossary}{}%
    {\vspace{-\intextsep}\centering
     \begin{longtable}{@{\hskip 1cm}p{3.5cm}p{\glsdescwidth}}}%
    {\end{longtable}}%
}
    \setabbreviationstyle [ acronym ] { long-short }
  }
  {
    \RequirePackage { glossaries }
    \setacronymstyle { short-long }
  }
\RenewDocumentCommand \genacrfullformat { m m } { \glsentrylong {#1} }
\bool_if:NTF \g__hit_en_bool { \makenoidxglossaries } { \makeglossaries }
%</thesis-glossary>
%    \end{macrocode}
%
% \Finale
